\section*{Abstract}

In enterprise analysis, STOP or No-Go outcomes are often interpreted as indicators of analytical failure, insufficient data, or lack of decisiveness. This paper challenges that interpretation by examining STOP within a formal framework in which enterprises are modeled as continua (K\_EA). The central question addressed is whether STOP can constitute a formally valid diagnostic outcome rather than a failure of analysis.

The analysis shifts the focus from decision preference to model admissibility. Analytical outcomes are treated as statements about what is representable within a given enterprise model under its internal constraints. Continuation is not assumed as a default result of competent analysis, but as a candidate outcome whose validity depends on the existence of admissible continuation states.

STOP is defined as the diagnostic outcome that arises when, relative to a specified enterprise model, no admissible continuation exists. Under this definition, STOP is a property of the modeled system rather than of the analyst or the analytical process. The paper states key propositions characterizing STOP, outlines their logical support, and specifies conditions under which the framework would be falsified.

The scope of the analysis is intentionally limited. It does not prescribe organizational actions, evaluate governance choices, or assess empirical prevalence. Its contribution is formal: it establishes that STOP can be a logically coherent and informative analytical outcome when admissibility of continuation is exhausted. By clarifying the status of STOP at the model level, the paper provides a foundation for subsequent methodological or practical discussions while remaining neutral with respect to their conclusions.
